% Generated by scripts/build_reports.py; edit reviewed Markdown inputs.
% Compile with XeLaTeX.
% Options for packages loaded elsewhere
\PassOptionsToPackage{unicode}{hyperref}
\PassOptionsToPackage{hyphens}{url}
\PassOptionsToPackage{dvipsnames,svgnames,x11names}{xcolor}
\documentclass[
  11pt,
]{article}
\usepackage{xcolor}
\usepackage[margin=1in]{geometry}
\usepackage{amsmath,amssymb}
\setcounter{secnumdepth}{-\maxdimen} % remove section numbering
\usepackage{iftex}
\ifPDFTeX
  \usepackage[T1]{fontenc}
  \usepackage[utf8]{inputenc}
  \usepackage{textcomp} % provide euro and other symbols
\else % if luatex or xetex
  \usepackage{unicode-math} % this also loads fontspec
  \defaultfontfeatures{Scale=MatchLowercase}
  \defaultfontfeatures[\rmfamily]{Ligatures=TeX,Scale=1}
\fi
\usepackage{lmodern}
\ifPDFTeX\else
  % xetex/luatex font selection
\fi
% Use upquote if available, for straight quotes in verbatim environments
\IfFileExists{upquote.sty}{\usepackage{upquote}}{}
\IfFileExists{microtype.sty}{% use microtype if available
  \usepackage[]{microtype}
  \UseMicrotypeSet[protrusion]{basicmath} % disable protrusion for tt fonts
}{}
\makeatletter
\@ifundefined{KOMAClassName}{% if non-KOMA class
  \IfFileExists{parskip.sty}{%
    \usepackage{parskip}
  }{% else
    \setlength{\parindent}{0pt}
    \setlength{\parskip}{6pt plus 2pt minus 1pt}}
}{% if KOMA class
  \KOMAoptions{parskip=half}}
\makeatother
\usepackage{longtable,booktabs,array}
\usepackage{caption}
\captionsetup[table]{skip=6pt}
\newcounter{none} % for unnumbered tables
\usepackage{calc} % for calculating minipage widths
% Correct order of tables after \paragraph or \subparagraph
\usepackage{etoolbox}
\makeatletter
\patchcmd\longtable{\par}{\if@noskipsec\mbox{}\fi\par}{}{}
\makeatother
% Allow footnotes in longtable head/foot
\IfFileExists{footnotehyper.sty}{\usepackage{footnotehyper}}{\usepackage{footnote}}
\makesavenoteenv{longtable}
\setlength{\emergencystretch}{3em} % prevent overfull lines
\providecommand{\tightlist}{%
  \setlength{\itemsep}{0pt}\setlength{\parskip}{0pt}}
\usepackage{amsmath,amssymb}
\usepackage{xurl}
\usepackage{fancyhdr}
\pagestyle{fancy}
\fancyhf{}
\fancyhead[L]{\small Pure Thought | Revised research specifications}
\fancyfoot[L]{\small 2026-09-17}
\fancyfoot[R]{\small\thepage}
\setlength{\headheight}{15pt}
\setlength{\emergencystretch}{3em}
\allowdisplaybreaks
\usepackage{bookmark}
\IfFileExists{xurl.sty}{\usepackage{xurl}}{} % add URL line breaks if available
\urlstyle{same}
% fallback for those not using the hyperref driver hyperxmp:
\makeatletter
\@ifundefined{xmpquote}{\newcommand{\xmpquote}[1]{#1}}{}
\makeatother
\hypersetup{
  pdftitle={14. Certified Weyl nodes and slice topology in a lattice model},
  pdfauthor={Pure Thought research portfolio},
  colorlinks=true,
  linkcolor={blue},
  filecolor={Maroon},
  citecolor={Blue},
  urlcolor={blue},
  pdfcreator={LaTeX via pandoc}}

\title{14. Certified Weyl nodes and slice topology in a lattice model}
\author{Pure Thought research portfolio}
\date{Revised 17 September 2026}

\begin{document}
\maketitle
\thispagestyle{empty}
\clearpage

\section{14. Certified Weyl nodes and slice topology in a lattice model}

\textbf{Revised specification --- 2026-09-11.}
\href{https://github.com/rick-stevens-ai/Pure-Thought/blob/main/reviews/14-Topological-Semimetals-Weyl-Dirac.md}{Detailed
critique} \ensuremath{\cdot}
\href{https://github.com/rick-stevens-ai/Pure-Thought/blob/main/archive/original-PRDs/14-Topological-Semimetals-Weyl-Dirac.md}{Original
description} \ensuremath{\cdot}
\href{https://github.com/rick-stevens-ai/Pure-Thought/blob/main/reviews/EVIDENCE-STANDARD.md}{Evidence
standard}

\subsection{Plain-language guide}

\subsubsection{The problem in everyday terms}

Energy bands in a three-dimensional crystal can meet at isolated points
instead of remaining separated everywhere. Near certain crossings, the
equations resemble those of a Weyl particle. Each crossing carries a
handedness-like topological label. This problem asks whether we can find
every such crossing in one lattice model, prove their labels, and
connect them to the behavior of two-dimensional momentum slices.

\subsubsection{Key terms}

\begin{itemize}
\item
  \textbf{Semimetal} --- A system in which the occupied and unoccupied
  electronic states are not separated by a full insulating gap; the
  selected model has isolated band crossings.
\item
  \textbf{Weyl node} --- A simple crossing of two bands with a
  characteristic linear structure and nonzero topological charge.
\item
  \textbf{Chirality} --- The signed handedness of a simple Weyl
  crossing, defined here from its velocity matrix.
\item
  \textbf{Velocity matrix} --- The derivatives describing how the
  Hamiltonian changes in each momentum direction near a node.
\item
  \textbf{Berry flux} --- The integrated Berry curvature through a
  closed surface around a crossing.
\item
  \textbf{Momentum slice} --- A two-dimensional section of the
  three-dimensional Brillouin zone, obtained by holding one momentum
  coordinate fixed.
\item
  \textbf{Fermi arc} --- An open surface-state contour whose relation to
  bulk nodes depends on the chosen surface and energy.
\item
  \textbf{Root isolation} --- A mathematical procedure proving that a
  small region contains a solution, often uniquely.
\item
  \textbf{Completeness} --- A proof that no additional nodes have been
  missed anywhere in the claimed domain.
\end{itemize}

\subsubsection{Why this matters}

Weyl nodes connect a local band crossing to a global topological
constraint. They help explain unusual electronic structures and surface
states, and provide controlled settings for studying the relation
between particle-like excitations and lattice physics. Missing even a
pair of nodes can change the interpretation of a phase diagram.

\subsubsection{What progress would mean}

A certified catalog would establish all crossings and slice invariants
of the selected model. It would give a reliable starting point for
surface and response calculations, without pretending that bulk nodes
uniquely determine every surface arc.

\subsection{Core question}

Can all nodes in a specified two-band three-dimensional Hamiltonian be
isolated, assigned chirality, and related to the Chern numbers of gapped
momentum slices?

\subsection{Scope and assumptions}

Begin with
\(H(k)=\sin k_x\,\sigma_x+\sin k_y\,\sigma_y+(2-\cos k_x-\cos k_y-\cos k_z)\sigma_z\),
k\ensuremath{\in}T\ensuremath{{}^{3}}, lower band occupied away from
nodes. Fix orientation and Pauli-matrix conventions. This exact
benchmark has nodes at (0,0,\ensuremath{\pm}\ensuremath{\pi}/2).

Subsequent searches must fix orbitals, hopping support, parameter box,
and symmetry representation. Dirac points and nodal lines are separate
extensions because their protection criteria differ.

\subsection{Mathematical target}

For H=d\ensuremath{{}_{0}}I+d\ensuremath{\cdot}\ensuremath{\sigma},
nodes solve d(k)=0. A simple node has invertible velocity matrix
\(V_{ij}=\partial d_i/\partial k_j\). Define Weyl chirality
\ensuremath{\chi}=sgn det V, and explicitly relate this to the
lower-band Chern flux under the chosen convention. Prove the complement
of all isolating neighborhoods has \textbar d(k)\textbar\textgreater0.

For fixed k\_z away from nodes, certify the two-dimensional occupied
Chern number. Its jumps track node charge. A slab calculation requires a
chosen surface normal and termination; arc details are not determined
solely by bulk node positions.

\subsection{Required outputs and proof obligations}

Return all root enclosures, uniqueness/completeness proof, Jacobian
signs, slice Chern numbers, and optional slab spectrum with error
control. Nodes need not have closed-form coordinates: certified
isolating boxes are acceptable. Robustness is scoped to perturbations
that prevent annihilation with opposite total charge.

\subsection{Validation and rejection controls}

Verify the two benchmark roots analytically and opposite determinant
signs. Check total charge on the periodic Brillouin zone vanishes.
Confirm slice invariants only on gapped slices. Test a perturbation
moving nodes and a parameter change permitting pair annihilation.

\subsection{Milestones and research extension}

Baseline: complete benchmark node and slice audit. Strong: certified
parameter-dependent node motion or one bounded model family. Research
extension: a symmetry-protected Dirac model with explicit little-group
constraints, or nodal-link invariants with their own definitions.

Milestones are acceptance gates, not calendar promises. Before a run,
freeze the concrete instance/domain, tolerance, compute budget, and
permitted dependencies. A bounded search that fails to find a result
must preserve an unresolved outcome.

\subsection{Starting instruction}

Work on the baseline above. First make every input and convention
explicit, then validate the reference case. Build the checker alongside
the calculation. Label numerical evidence, finite-model certificates,
and continuum theorems separately; return the strongest justified result
and identify the exact unmet proof obligation.

\subsection{Primary references and their role}

\begin{itemize}
\tightlist
\item
  \href{https://arxiv.org/abs/1007.0016}{Wan et al., Electronic
  Structure of Pyrochlore Iridates (published as Topological semimetal
  and Fermi-arc surface states)} --- Reference for Weyl nodes and
  surface arcs; benchmark and proof obligations here are explicitly
  specified.
\end{itemize}

References support definitions or limitations; the scope choices and
proposed milestones are editorial recommendations, not claims that these
results have already been obtained.

\subsection{Displayed formulation and claim boundary}

\href{https://github.com/rick-stevens-ai/Pure-Thought/blob/main/PRDs/14-Topological-Semimetals-Weyl-Dirac.md}{Canonical
specification}.

\[d(k)=(\sin k_x,\sin k_y,2-\cos k_x-\cos k_y-\cos k_z),
\quad d(k_*)=0,\]

\[k_*=(0,0,\pm\pi/2),\qquad
V_{ij}=\partial_{k_j}d_i(k_*),\qquad \chi=\operatorname{sgn}\det V.\]

The two determinants are \(+1\) and \(-1\) in the stated coordinate
order. Relate this convention explicitly to the occupied-band flux.
Certify absence of other nodes on the complement and the topology of
gapped slices; surface arcs additionally depend on the chosen boundary.

\subsection{Detailed review and rationale}

\href{https://github.com/rick-stevens-ai/Pure-Thought/blob/main/PRDs/14-Topological-Semimetals-Weyl-Dirac.md}{Revised
problem} \ensuremath{\cdot}
\href{https://github.com/rick-stevens-ai/Pure-Thought/blob/main/archive/original-PRDs/14-Topological-Semimetals-Weyl-Dirac.md}{Preserved
original} \ensuremath{\cdot}
\href{https://github.com/rick-stevens-ai/Pure-Thought/blob/f854d4e12b55c1f5b5c2a964b75f2d23496d981e/PRDs/14-Topological-Semimetals-Weyl-Dirac.md}{Pinned
upstream source}

Line references below refer to the preserved original at commit
\texttt{f854d4e12b55}, not the rewritten file. ``Blocking'' means the
original issue can invalidate the scientific target or claimed
certificate; ``major'' means it materially changes scope, inference, or
acceptance criteria. These are critiques of specifications, not results
of running the embedded research-code sketches.

\subsubsection{Detailed findings}

\begin{enumerate}
\def\labelenumi{\arabic{enumi}.}
\item
  \textbf{Major --- oversimplified symmetry (22--25).} Time reversal and
  inversion alone are not a universal recipe for a protected fourfold
  Dirac crossing; representation and additional symmetry constraints
  matter. They cannot replace an explicit model.
\item
  \textbf{Major --- isotropy assumed (73--87).} A general Weyl cone has
  a velocity matrix, can be tilted, and is not specified by one scalar
  v\_F. Its determinant gives the simple-node chirality in a declared
  convention.
\item
  \textbf{Major --- completeness (99--105).} Finding roots and checking
  their charge sum does not prove there are no additional zero-charge
  pairs. A proof on the complement of the isolating boxes is needed.
\item
  \textbf{Major --- boundary overclaim.} Fermi-arc connectivity depends
  on surface projection and termination; opposite nodes can project
  together. The revision uses slice topology as the robust intermediate
  claim and separates surface-specific predictions.
\end{enumerate}

\subsubsection{What the revision preserves and changes}

The central motivation is retained, but the executable question is now:
\textbf{Can all nodes in a specified two-band three-dimensional
Hamiltonian be isolated, assigned chirality, and related to the Chern
numbers of gapped momentum slices?} The revised scope chooses a concrete
starting model; those choices are proposed research-design decisions,
not facts inferred from the original. The baseline can succeed without
resolving the full open research program.

The replacement also removes the original implementation sketches from
the active specification. They remain in the archive for comparison, but
have not been made executable or certified. New implementations should
derive their equations from the revised definitions rather than copy
unchecked prototypes.

\subsubsection{Acceptance criteria to use instead}

Baseline: complete benchmark node and slice audit. Strong: certified
parameter-dependent node motion or one bounded model family. Research
extension: a symmetry-protected Dirac model with explicit little-group
constraints, or nodal-link invariants with their own definitions.

\subsubsection{Sources supporting the corrections}

\begin{itemize}
\tightlist
\item
  \href{https://arxiv.org/abs/1007.0016}{Wan et al., Electronic
  Structure of Pyrochlore Iridates (published as Topological semimetal
  and Fermi-arc surface states)} --- Reference for Weyl nodes and
  surface arcs; benchmark and proof obligations here are explicitly
  specified.
\end{itemize}

\section{Evidence and execution standard}

This standard applies to the revised specifications. It is a proposed
contract for future research runs; this editorial revision does not
supply the scientific certificates themselves.

\subsection{What ``pure thought'' means here}

Use explicitly stated mathematical models and deductions without fitting
conclusions to experimental outcomes. Synthetic examples, exact model
inputs, analytic references, numerical search, and established software
are allowed. Physical constants, electronic-structure approximations,
constitutive laws, and empirical parameter sets must be identified when
used. A model-derived result is conditional on that model.

The original ban on legacy software conflicts with its proposed reliance
on NumPy, SciPy, CVXPY, nauty, and other libraries. Reimplementing every
optimizer and algebra system is not necessary for scientific
independence. Prefer a small independently checkable certificate and a
documented trusted base. A fresh implementation can be a project
objective, but must not be confused with a mathematical proof or novelty
claim.

\subsection{Four evidence levels}

{\def\LTcaptype{none} % do not increment counter
\begin{longtable}[]{@{}
  >{\raggedright\arraybackslash}p{(\linewidth - 4\tabcolsep) * \real{0.3333}}
  >{\raggedright\arraybackslash}p{(\linewidth - 4\tabcolsep) * \real{0.3333}}
  >{\raggedright\arraybackslash}p{(\linewidth - 4\tabcolsep) * \real{0.3333}}@{}}
\toprule\noalign{}
\begin{minipage}[b]{\linewidth}\raggedright
Label
\end{minipage} & \begin{minipage}[b]{\linewidth}\raggedright
What it establishes
\end{minipage} & \begin{minipage}[b]{\linewidth}\raggedright
What it does not establish
\end{minipage} \\
\midrule\noalign{}
\endhead
\bottomrule\noalign{}
\endlastfoot
Numerical evidence & A reproducible computation at stated precision,
sample size, and discretization & Exact truth, continuum coverage, or a
universal theorem \\
Finite-model certificate & An exact/validated statement about the
specified finite model or finite relaxation & An omitted tail,
thermodynamic limit, or physical realization \\
Conditional mathematical theorem & A proof over the stated domain with
every hypothesis and error controlled & Truth outside those assumptions
or a demonstrated physical system \\
Formal proof & A checked theorem in a named proof assistant, with
imported axioms/dependencies listed & Correct physical modeling merely
by virtue of formalization \\
\end{longtable}
}

Formalization is an optional extra deliverable unless the specific
project makes it mandatory. High precision, JSON output, solver status
and small residuals do not by themselves confer any of these proof
levels.

\subsection{Claim directions}

For minimization over a larger feasible set, the optimum is a lower
bound on the original minimum. A physical feasible point supplies an
upper bound. For maximization the directions reverse. Track feasibility
of the primal/dual witnesses separately from their objective values.

Infeasibility of an inner ansatz excludes only that ansatz.
Infeasibility of an outer relaxation can exclude the original problem if
the map into the relaxation is proved and all approximations preserve
necessity. Passing necessary tests means ``not excluded''; existence
needs a constructive sufficient argument.

A smallest observed example is not a minimality theorem. A complete
catalog must define the universe, equivalence relation, boundaries, and
proof that pruning misses nothing. A finite-size numerical trend does
not prove an asymptotic limit.

\subsection{Minimum run manifest}

Before calculation, record the problem ID and revision hash, exact
inputs and conventions, quantified domain, equivalence relation, target
observable, accepted evidence level, tolerances, and resource limits.
Record dependencies and precision/rounding mode. Where randomization is
used, preserve seeds, trial counts and raw sufficient statistics.
Resource limits are planning decisions, not evidence that a mathematical
solution cannot exist.

An output bundle should include:

\begin{itemize}
\tightlist
\item
  The statement being claimed, including all assumptions and
  quantifiers.
\item
  Construction or proof object, plus an independently runnable verifier
  and its trust assumptions.
\item
  Numerical/error ledger separating roundoff, discretization,
  truncation, optimization, sampling and modeling error where relevant.
\item
  Positive reference cases and negative controls that the verifier must
  reject.
\item
  Coverage: included domain, omitted domain, unresolved regions and
  termination reason.
\end{itemize}

These are requirements for future implementations, not a promise of
polynomial-size certificates for every problem.

\subsection{Checking specific certificate types}

\textbf{LP/SDP/SOS:} verify equality constraints, objective direction,
and cone feasibility using exact arithmetic or outward-rounded
intervals. Rationalizing a numerical solution can break PSD constraints.
Correct residuals with justified norm bounds; do not silently project
and reuse the old objective. Distinguish a proof of a polynomial
approximation from a proof for the original function.

\textbf{Interval numerics:} state function enclosures, domains and
outward rounding. A small root box needs existence and, if claimed,
uniqueness. Cover the complement for a global enumeration. Include
singularities and tails rather than discarding them numerically.

\textbf{SAT/SMT:} a solver's UNSAT label is not a certificate. Preserve
a proof supported by a checker and the proof that the encoding, symmetry
reduction and blocked solutions match the mathematical universe. DRAT
applies to an appropriate Boolean CNF proof, not automatically to every
real-arithmetic or SDP claim.

\textbf{Statistics:} specify the estimand, dependence assumptions,
confidence procedure and bias. Zero observed failures does not prove
zero failure probability. MCMC error bars based on estimated
autocorrelation are not generally rigorous mixing guarantees.

\subsection{Milestones and novelty}

Baseline means a working, independently validated result on a concrete
reference problem. Strong means a broader or sharper result with its
claimed evidence level. Research extension means a possible novel result
whose feasibility and significance need review. A negative theorem can
be valuable; an unsuccessful search can also be useful if its scope is
documented, but the two are not interchangeable.

Do not promise ``publication-quality'' in a fixed number of months.
Before claiming novelty, conduct a focused current literature comparison
against the exact theorem/model/metric. The sources in this revision
check key definitions and limitations; they are not an exhaustive
novelty audit across all 30 fields.

\end{document}
